\documentclass[10pt,a4paper]{amsart}
\usepackage{amsmath,amssymb,amsthm,graphicx}
% Retained preamble for the Stacks Project chapter morphisms.tex (Morphisms of Schemes).
% The chapter file pinned-original-source.tex carries only \input{preamble} on its first line, so no
% preamble text lives inside the chapter itself.  The Stacks Project uses one single shared preamble for
% every chapter; the snapshot stages it verbatim in sources/stacks-chow/preamble.tex.  The block below is
% copied mechanically, line for line, from that file's lines 14-18 (the Xy-pic setup) and lines 154-217
% (the native theorem declarations together with the mathematical macros).  Nothing was added, removed or
% redefined.  The wrapper class and the amsmath/amssymb/amsthm/graphicx base packages are supplied by the
% assembler, so the class/\externaldocument part of the shared preamble is not repeated here.
% For commutative diagrams we use Xy-pic
\usepackage[all]{xy}

% We use 2cell for 2-commutative diagrams.
\xyoption{2cell}
% Theorem environments.
%
\theoremstyle{plain}
\newtheorem{theorem}[subsection]{Theorem}
\newtheorem{proposition}[subsection]{Proposition}
\newtheorem{lemma}[subsection]{Lemma}

\theoremstyle{definition}
\newtheorem{definition}[subsection]{Definition}
\newtheorem{example}[subsection]{Example}
\newtheorem{exercise}[subsection]{Exercise}
\newtheorem{situation}[subsection]{Situation}

\theoremstyle{remark}
\newtheorem{remark}[subsection]{Remark}
\newtheorem{remarks}[subsection]{Remarks}

\numberwithin{equation}{subsection}

% Macros
%
\def\lim{\mathop{\mathrm{lim}}\nolimits}
\def\colim{\mathop{\mathrm{colim}}\nolimits}
\def\Spec{\mathop{\mathrm{Spec}}}
\def\Hom{\mathop{\mathrm{Hom}}\nolimits}
\def\Ext{\mathop{\mathrm{Ext}}\nolimits}
\def\SheafHom{\mathop{\mathcal{H}\!\mathit{om}}\nolimits}
\def\SheafExt{\mathop{\mathcal{E}\!\mathit{xt}}\nolimits}
\def\Sch{\mathit{Sch}}
\def\Mor{\mathop{\mathrm{Mor}}\nolimits}
\def\Ob{\mathop{\mathrm{Ob}}\nolimits}
\def\Sh{\mathop{\mathit{Sh}}\nolimits}
\def\NL{\mathop{N\!L}\nolimits}
\def\CH{\mathop{\mathrm{CH}}\nolimits}
\def\proetale{{pro\text{-}\acute{e}tale}}
\def\etale{{\acute{e}tale}}
\def\QCoh{\mathit{QCoh}}
\def\Ker{\mathop{\mathrm{Ker}}}
\def\Im{\mathop{\mathrm{Im}}}
\def\Coker{\mathop{\mathrm{Coker}}}
\def\Coim{\mathop{\mathrm{Coim}}}

% Boxtimes
%
\DeclareMathSymbol{\boxtimes}{\mathbin}{AMSa}{"02}

%
% Macros for moduli stacks/spaces
%
\def\QCohstack{\mathcal{QC}\!\mathit{oh}}
\def\Cohstack{\mathcal{C}\!\mathit{oh}}
\def\Spacesstack{\mathcal{S}\!\mathit{paces}}
\def\Quotfunctor{\mathrm{Quot}}
\def\Hilbfunctor{\mathrm{Hilb}}
\def\Curvesstack{\mathcal{C}\!\mathit{urves}}
\def\Polarizedstack{\mathcal{P}\!\mathit{olarized}}
\def\Complexesstack{\mathcal{C}\!\mathit{omplexes}}
% \Pic is the operator that assigns to X its picard group, usage \Pic(X)
% \Picardstack_{X/B} denotes the Picard stack of X over B
% \Picardfunctor_{X/B} denotes the Picard functor of X over B
\def\Pic{\mathop{\mathrm{Pic}}\nolimits}
\def\Picardstack{\mathcal{P}\!\mathit{ic}}
\def\Picardfunctor{\mathrm{Pic}}
\def\Deformationcategory{\mathcal{D}\!\mathit{ef}}
\usepackage{cleveref}
\crefname{originalreference}{original reference}{original references}
\Crefname{originalreference}{Original reference}{Original references}
\makeatletter
\newlabel{algebra-equation-good-locus}{{algebra-equation-good-locus}{source}{}{sourcepointer-1}{}}
\newlabel{algebra-equation-good-locus@cref}{{[originalreference][0][]algebra-equation-good-locus}{[1][0][]source}{}{}{}}
\makeatother
\crefname{originalreference}{original reference}{original references}
\Crefname{originalreference}{Original reference}{Original references}
\makeatletter
\newlabel{algebra-lemma-generic-flatness}{{algebra-lemma-generic-flatness}{source}{}{sourcepointer-2}{}}
\newlabel{algebra-lemma-generic-flatness@cref}{{[originalreference][0][]algebra-lemma-generic-flatness}{[1][0][]source}{}{}{}}
\makeatother
\crefname{originalreference}{original reference}{original references}
\Crefname{originalreference}{Original reference}{Original references}
\makeatletter
\newlabel{algebra-lemma-generic-flatness-reduced}{{algebra-lemma-generic-flatness-reduced}{source}{}{sourcepointer-3}{}}
\newlabel{algebra-lemma-generic-flatness-reduced@cref}{{[originalreference][0][]algebra-lemma-generic-flatness-reduced}{[1][0][]source}{}{}{}}
\makeatother
\crefname{originalreference}{original reference}{original references}
\Crefname{originalreference}{Original reference}{Original references}
\makeatletter
\newlabel{algebra-lemma-support-finite-presentation-constructible}{{algebra-lemma-support-finite-presentation-constructible}{source}{}{sourcepointer-4}{}}
\newlabel{algebra-lemma-support-finite-presentation-constructible@cref}{{[originalreference][0][]algebra-lemma-support-finite-presentation-constructible}{[1][0][]source}{}{}{}}
\makeatother
\crefname{originallemma}{original lemma}{original lemmas}
\Crefname{originallemma}{Original lemma}{Original lemmas}
\makeatletter
\newlabel{lemma-flat-module-characterize}{{lemma-flat-module-characterize}{source}{}{sourcepointer-5}{}}
\newlabel{lemma-flat-module-characterize@cref}{{[originallemma][0][]lemma-flat-module-characterize}{[1][0][]source}{}{}{}}
\makeatother
\crefname{originallemma}{original lemma}{original lemmas}
\Crefname{originallemma}{Original lemma}{Original lemmas}
\makeatletter
\newlabel{lemma-locally-finite-presentation-characterize}{{lemma-locally-finite-presentation-characterize}{source}{}{sourcepointer-6}{}}
\newlabel{lemma-locally-finite-presentation-characterize@cref}{{[originallemma][0][]lemma-locally-finite-presentation-characterize}{[1][0][]source}{}{}{}}
\makeatother
\crefname{originallemma}{original lemma}{original lemmas}
\Crefname{originallemma}{Original lemma}{Original lemmas}
\makeatletter
\newlabel{lemma-locally-finite-type-characterize}{{lemma-locally-finite-type-characterize}{source}{}{sourcepointer-7}{}}
\newlabel{lemma-locally-finite-type-characterize@cref}{{[originallemma][0][]lemma-locally-finite-type-characterize}{[1][0][]source}{}{}{}}
\makeatother
\crefname{originalreference}{original reference}{original references}
\Crefname{originalreference}{Original reference}{Original references}
\makeatletter
\newlabel{properties-lemma-characterize-reduced}{{properties-lemma-characterize-reduced}{source}{}{sourcepointer-8}{}}
\newlabel{properties-lemma-characterize-reduced@cref}{{[originalreference][0][]properties-lemma-characterize-reduced}{[1][0][]source}{}{}{}}
\makeatother
\crefname{originalreference}{original reference}{original references}
\Crefname{originalreference}{Original reference}{Original references}
\makeatletter
\newlabel{properties-lemma-finite-presentation-module}{{properties-lemma-finite-presentation-module}{source}{}{sourcepointer-9}{}}
\newlabel{properties-lemma-finite-presentation-module@cref}{{[originalreference][0][]properties-lemma-finite-presentation-module}{[1][0][]source}{}{}{}}
\makeatother
\crefname{originalreference}{original reference}{original references}
\Crefname{originalreference}{Original reference}{Original references}
\makeatletter
\newlabel{properties-lemma-finite-type-module}{{properties-lemma-finite-type-module}{source}{}{sourcepointer-10}{}}
\newlabel{properties-lemma-finite-type-module@cref}{{[originalreference][0][]properties-lemma-finite-type-module}{[1][0][]source}{}{}{}}
\makeatother
\crefname{originalproposition}{original proposition}{original propositions}
\Crefname{originalproposition}{Original proposition}{Original propositions}
\makeatletter
\newlabel{proposition-generic-flatness}{{proposition-generic-flatness}{source}{}{sourcepointer-11}{}}
\newlabel{proposition-generic-flatness@cref}{{[originalproposition][0][]proposition-generic-flatness}{[1][0][]source}{}{}{}}
\makeatother
\crefname{originalreference}{original reference}{original references}
\Crefname{originalreference}{Original reference}{Original references}
\makeatletter
\newlabel{schemes-lemma-characterize-quasi-separated}{{schemes-lemma-characterize-quasi-separated}{source}{}{sourcepointer-12}{}}
\newlabel{schemes-lemma-characterize-quasi-separated@cref}{{[originalreference][0][]schemes-lemma-characterize-quasi-separated}{[1][0][]source}{}{}{}}
\makeatother
\title{Generic flatness over a reduced base: an open dense locus on which a finite type morphism and a finite type sheaf are flat and of finite presentation}
\author{The Stacks Project Authors}
\date{stacks-project:morphisms.tex:local-snapshot-20261012; GNU FDL 1.2 or later}
\begin{document}
\maketitle
\noindent\textit{Editorial note.} All delivered mathematical text is the author's own, copied from the pinned Stacks Project chapter file pinned-original-source.tex; no statement or proof step was rewritten, reordered or removed.\par\medskip
\noindent\textit{Editorial note.} Following the multi-claim convention for this source, source\_url is the source-level URL https://stacks.math.columbia.edu/ and the source label proposition-generic-flatness-reduced is carried in claim\_key and here rather than in a per-claim tag URL.\par\medskip
\noindent\textit{Editorial note.} The preamble is the chapter's shared Stacks Project preamble, retained as sources/stacks-morphisms/retained-preamble.tex: a line-for-line excerpt of the shared Stacks preamble staged with this snapshot. No macro was invented or redefined, and this card adds no preamble text of its own.\par\medskip
\noindent\textit{Editorial note.} Statement span 5246-5260 and proof span 5262-5345 come from the halfA pool card operations/pool-cards/stacks-morphisms-halfA.json and were re-read in the pinned chapter file before this card was written.\par\medskip
\noindent\textit{Editorial note.} The counted proof opens with the author's sentence 'For the impatient reader: This proof is a repeat of the proof of Proposition [generic flatness] using Algebra, Lemma [generic flatness, reduced] instead of Algebra, Lemma [generic flatness].' That sentence is kept verbatim; the rest of the span is the actual proof, not a pointer.\par\medskip
\noindent\textit{Editorial note.} The algebra input is cited by label only and is delivered as an external navigation pointer; no literal text of the cited Algebra lemma was copied in, and no content of it is asserted from memory.\par\medskip
\noindent\textit{Editorial note.} 12 cross-references occur inside the counted spans; they are listed in external\_references as navigation pointers only.\par\medskip
\noindent\textit{Allowed standard prerequisites (editorial).}\begin{itemize}
\item The algebraic generic flatness input for a reduced base ring, invoked by the author as Algebra, Lemma algebra-lemma-generic-flatness-reduced and referenced from the counted span; its statement is not copied into this card and no content is asserted from memory. Applied at proof lines 5287-5289 to the pairs (A to B\_i, M\_ij) and (A to B\_i, M\_i). Supplies the dense opens that are intersected to build the announced open subscheme U.
\item Locality on the base of flatness and of finite presentation, as characterised in the same chapter by the two lemmas the author cites. Cited at proof lines 5264-5267. Lets the argument assume that the base is affine with reduced coordinate ring.
\item The affine dictionary used repeatedly in the proof: finite type morphisms correspond to finite type algebras, finite type quasi-coherent modules on an affine scheme correspond to finite modules, and finite presentation of a module is detected on the coordinate ring. Used at proof lines 5274-5280 and again at 5338-5341. Translates the geometric hypotheses into the algebra input and translates the algebra output back into flatness and finite presentation of X over S and of F over O\_X.
\end{itemize}\medskip
\ifdefined\hypertarget\hypertarget{sourcepointer-1}{}\fi
\noindent\textit{Source reference.} Original reference algebra-equation-good-locus: sources/stacks-morphisms/pinned-original-source.tex, source lines 5300--5300. Same-chapter reference used in the retained author text at source line(s) 5300. It is supplied only as a navigation pointer so that the retained text compiles; the counted item delivers the author's own statement and the author's complete proof, and no part of the referenced item is copied in or substituted for the counted argument.\par\medskip
\ifdefined\hypertarget\hypertarget{sourcepointer-2}{}\fi
\noindent\textit{Source reference.} Original reference algebra-lemma-generic-flatness: sources/stacks-morphisms/pinned-original-source.tex, source lines 5210--5210. Same-chapter reference used in the retained author text at source line(s) 5210, 5268. It is supplied only as a navigation pointer so that the retained text compiles; the counted item delivers the author's own statement and the author's complete proof, and no part of the referenced item is copied in or substituted for the counted argument.\par\medskip
\ifdefined\hypertarget\hypertarget{sourcepointer-3}{}\fi
\noindent\textit{Source reference.} Original reference algebra-lemma-generic-flatness-reduced: sources/stacks-morphisms/pinned-original-source.tex, source lines 5266--5266. Same-chapter reference used in the retained author text at source line(s) 5266, 5295. It is supplied only as a navigation pointer so that the retained text compiles; the counted item delivers the author's own statement and the author's complete proof, and no part of the referenced item is copied in or substituted for the counted argument.\par\medskip
\ifdefined\hypertarget\hypertarget{sourcepointer-4}{}\fi
\noindent\textit{Source reference.} Original reference algebra-lemma-support-finite-presentation-constructible: sources/stacks-morphisms/pinned-original-source.tex, source lines 5236--5236. Same-chapter reference used in the retained author text at source line(s) 5236, 5341. It is supplied only as a navigation pointer so that the retained text compiles; the counted item delivers the author's own statement and the author's complete proof, and no part of the referenced item is copied in or substituted for the counted argument.\par\medskip
\ifdefined\hypertarget\hypertarget{sourcepointer-5}{}\fi
\noindent\textit{Source reference.} Original lemma lemma-flat-module-characterize: sources/stacks-morphisms/pinned-original-source.tex, source lines 4660--4660. Same-chapter reference used in the retained author text at source line(s) 4717, 4730, 5273. It is supplied only as a navigation pointer so that the retained text compiles; the counted item delivers the author's own statement and the author's complete proof, and no part of the referenced item is copied in or substituted for the counted argument.\par\medskip
\ifdefined\hypertarget\hypertarget{sourcepointer-6}{}\fi
\noindent\textit{Source reference.} Original lemma lemma-locally-finite-presentation-characterize: sources/stacks-morphisms/pinned-original-source.tex, source lines 4096--4096. Same-chapter reference used in the retained author text at source line(s) 4141, 4166, 4288, 4818, 5230, 5274, 5335, 5525, 6717. It is supplied only as a navigation pointer so that the retained text compiles; the counted item delivers the author's own statement and the author's complete proof, and no part of the referenced item is copied in or substituted for the counted argument.\par\medskip
\ifdefined\hypertarget\hypertarget{sourcepointer-7}{}\fi
\noindent\textit{Source reference.} Original lemma lemma-locally-finite-type-characterize: sources/stacks-morphisms/pinned-original-source.tex, source lines 2673--2673. Same-chapter reference used in the retained author text at source line(s) 2717, 2739, 2801, 2863, 3019, 3035, 3157, 3348, 3613, 3707, 3868, 5197, 5282, 6702, 9325, 9390, 9440, 11421, 13422, 13686. It is supplied only as a navigation pointer so that the retained text compiles; the counted item delivers the author's own statement and the author's complete proof, and no part of the referenced item is copied in or substituted for the counted argument.\par\medskip
\ifdefined\hypertarget\hypertarget{sourcepointer-8}{}\fi
\noindent\textit{Source reference.} Original reference properties-lemma-characterize-reduced: sources/stacks-morphisms/pinned-original-source.tex, source lines 5277--5277. Same-chapter reference used in the retained author text at source line(s) 5277. It is supplied only as a navigation pointer so that the retained text compiles; the counted item delivers the author's own statement and the author's complete proof, and no part of the referenced item is copied in or substituted for the counted argument.\par\medskip
\ifdefined\hypertarget\hypertarget{sourcepointer-9}{}\fi
\noindent\textit{Source reference.} Original reference properties-lemma-finite-presentation-module: sources/stacks-morphisms/pinned-original-source.tex, source lines 5232--5232. Same-chapter reference used in the retained author text at source line(s) 5232, 5337, 6718. It is supplied only as a navigation pointer so that the retained text compiles; the counted item delivers the author's own statement and the author's complete proof, and no part of the referenced item is copied in or substituted for the counted argument.\par\medskip
\ifdefined\hypertarget\hypertarget{sourcepointer-10}{}\fi
\noindent\textit{Source reference.} Original reference properties-lemma-finite-type-module: sources/stacks-morphisms/pinned-original-source.tex, source lines 703--703. Same-chapter reference used in the retained author text at source line(s) 703, 754, 4224, 5201, 5286, 6703, 10562, 11247. It is supplied only as a navigation pointer so that the retained text compiles; the counted item delivers the author's own statement and the author's complete proof, and no part of the referenced item is copied in or substituted for the counted argument.\par\medskip
\ifdefined\hypertarget\hypertarget{sourcepointer-11}{}\fi
\noindent\textit{Source reference.} Original proposition proposition-generic-flatness: sources/stacks-morphisms/pinned-original-source.tex, source lines 5172--5172. Same-chapter reference used in the retained author text at source line(s) 5264, 13542. It is supplied only as a navigation pointer so that the retained text compiles; the counted item delivers the author's own statement and the author's complete proof, and no part of the referenced item is copied in or substituted for the counted argument.\par\medskip
\ifdefined\hypertarget\hypertarget{sourcepointer-12}{}\fi
\noindent\textit{Source reference.} Original reference schemes-lemma-characterize-quasi-separated: sources/stacks-morphisms/pinned-original-source.tex, source lines 1763--1763. Same-chapter reference used in the retained author text at source line(s) 1763, 2254, 5238, 5343. It is supplied only as a navigation pointer so that the retained text compiles; the counted item delivers the author's own statement and the author's complete proof, and no part of the referenced item is copied in or substituted for the counted argument.\par\medskip
\begin{proposition}[Generic flatness, reduced case]
\label{proposition-generic-flatness-reduced}
Let $f : X \to S$ be a morphism of schemes.
Let $\mathcal{F}$ be a quasi-coherent sheaf of $\mathcal{O}_X$-modules.
Assume
\begin{enumerate}
\item $S$ is reduced,
\item $f$ is of finite type, and
\item $\mathcal{F}$ is a finite type $\mathcal{O}_X$-module.
\end{enumerate}
Then there exists an open dense subscheme $U \subset S$ such that
$X_U \to U$ is flat and of finite presentation and such that
$\mathcal{F}|_{X_U}$ is flat over $U$ and of finite presentation
over $\mathcal{O}_{X_U}$.
\end{proposition}

\begin{proof}
For the impatient reader: This proof is a repeat of the proof of
Proposition \ref{proposition-generic-flatness}
using
Algebra, Lemma \ref{algebra-lemma-generic-flatness-reduced}
instead of
Algebra, Lemma \ref{algebra-lemma-generic-flatness}.

\medskip\noindent
Since being flat and being of finite presentation is local on the
base, see
Lemmas \ref{lemma-flat-module-characterize} and
\ref{lemma-locally-finite-presentation-characterize},
we may work affine locally on $S$. Thus we may assume that
$S = \Spec(A)$, where $A$ is a reduced ring (see
Properties, Lemma \ref{properties-lemma-characterize-reduced}).
As $f$ is of finite type, it is quasi-compact, so $X$ is quasi-compact.
Hence we can find a finite affine open cover
$X = \bigcup_{i = 1, \ldots, n} X_i$. Write $X_i = \Spec(B_i)$.
Then $B_i$ is a finite type $A$-algebra, see
Lemma \ref{lemma-locally-finite-type-characterize}.
Moreover there are finite type
$B_i$-modules $M_i$ such that $\mathcal{F}|_{X_i}$ is the
quasi-coherent sheaf associated to the $B_i$-module $M_i$, see
Properties, Lemma \ref{properties-lemma-finite-type-module}.
Next, for each pair of indices $i, j$ choose an ideal $I_{ij} \subset B_i$
such that $X_i \setminus X_i \cap X_j = V(I_{ij})$ inside
$X_i = \Spec(B_i)$. Set $M_{ij} = B_i/I_{ij}$ and think
of it as a $B_i$-module. Then $V(I_{ij}) = \text{Supp}(M_{ij})$
and $M_{ij}$ is a finite $B_i$-module.

\medskip\noindent
At this point we apply
Algebra, Lemma \ref{algebra-lemma-generic-flatness-reduced} to
the pairs $(A \to B_i, M_{ij})$ and to the pairs $(A \to B_i, M_i)$.
Thus we obtain dense opens
$U(A \to B_i, M_{ij}) \subset S$ and dense opens
$U(A \to B_i, M_i) \subset S$ with notation as in
Algebra, Equation (\ref{algebra-equation-good-locus}).
Since a finite intersection of dense opens is dense open, we see that
$$
U =
\bigcap\nolimits_{i, j} U(A \to B_i, M_{ij})
\quad\cap\quad
\bigcap\nolimits_i U(A \to B_i, M_i)
$$
is open and dense in $S$. We claim that $U$ is the desired open.

\medskip\noindent
Pick $u \in U$. By definition of the loci $U(A \to B_i, M_{ij})$
and $U(A \to B, M_i)$ there exist $f_{ij}, f_i \in A$ such that
(a) $u \in D(f_i)$ and $u \in D(f_{ij})$,
(b) $A_{f_{ij}} \to B_{i, f_{ij}}$ is flat and of finite presentation
and $M_{ij, f_{ij}}$ is flat over $A_{f_{ij}}$ and of finite presentation
over $B_{i, f_{ij}}$, and
(c) $B_{i, f_i}$ is flat and of finite presentation over $A_{f_i}$ and
$M_{i, f_i}$ is of finite presentation over $B_{i, f_i}$ and flat over
$A_{f_i}$. Set
$f = (\prod f_i) (\prod f_{ij})$.
Now it suffices to prove that $X \to S$ is flat and of finite presentation
over $D(f)$ and that $\mathcal{F}$ restricted to $X_{D(f)}$ is
flat over $D(f)$ and of finite presentation over the structure sheaf
of $X_{D(f)}$.

\medskip\noindent
Hence we may replace $A$ by $A_f$, i.e.,
perform the base change by $\Spec(A_f) \to S$.
After this base change we see that each of $A \to B_i$ is
flat and of finite presentation and that $M_i$, $M_{ij}$ are flat over $A$
and of finite presentation over $B_i$.
This already proves that $X \to S$ is quasi-compact,
locally of finite presentation, flat, and that $\mathcal{F}$
is flat over $S$ and of finite presentation over $\mathcal{O}_X$, see
Lemma \ref{lemma-locally-finite-presentation-characterize}
and
Properties, Lemma \ref{properties-lemma-finite-presentation-module}.
Since $M_{ij}$ is of finite presentation over $B_i$ we see that
$X_i \cap X_j = X_i \setminus \text{Supp}(M_{ij})$ is a quasi-compact
open of $X_i$, see
Algebra, Lemma \ref{algebra-lemma-support-finite-presentation-constructible}.
Hence we see that $X \to S$ is quasi-separated by
Schemes, Lemma \ref{schemes-lemma-characterize-quasi-separated}.
This proves the proposition.
\end{proof}

\end{document}
